\subsection{Do the diagnostics detect punishment where it exists?}
\label{sec:validation}

A test that always says ``no punishment'' would be useless. Before applying
the diagnostics to traders, we run them on the case where algorithmic
collusion is held to be sustained by punishment: the logit Bertrand duopoly
of \citet{calvano2020}, with their baseline parameters (demand
$a_i=2$, $a_0=0$, $\mu=1/4$, cost $1$, 15 prices spanning the Nash-to-monopoly
range plus 10\% on each side, memory of both last prices, $\alpha=0.15$,
$\beta=4\times10^{-6}$, $\delta=0.95$; 100 sessions, $2.5\times10^6$ periods).
The deviation test is the one used for traders: one firm plays its static
best response to the rival's greedy price for one period, paired with an
undeviated copy. Table~\ref{tab:validation} reports the profit-based collusion
index, the rival's price change one period after the deviation (in units of
the monopoly--Nash price gap), and the deviator's discounted gain over 16
periods (in units of the per-period monopoly--Nash profit gap).

\begin{table}[t]
  \centering
  \caption{Diagnostics outside the Kyle market. Logit Bertrand: the
  \citet{calvano2020} baseline and placebos. Dealers: two Q-learning market
  makers quoting asks against informed and liquidity buyers
  (Section~\ref{sec:dealers}). Rival response: price change one period after
  a forced best-response deviation, relative to the monopoly--Nash price gap.
  Deviator gain: discounted over 16 periods, relative to the per-period
  monopoly--Nash profit gap. Verdict ``punishment'' when the rival cuts its
  price and deviating loses, both significantly. 100 sessions; 95\% CIs.}
  \label{tab:validation}
  \resizebox{\linewidth}{!}{\begin{tabular}{lccccl}
\toprule
Learners & $\Delta$ profit & $\Delta$ best ask & Rival price response & Deviator gain & Verdict \\
\midrule
\multicolumn{6}{l}{\emph{Logit Bertrand (Calvano et al.)}} \\
\quad Baseline ($\gamma=0.95$, memory) & $0.86 \pm 0.02$ & -- & $-0.54 \pm 0.06$ & $-1.72 \pm 0.22$ & punishment \\
\quad Myopic ($\gamma=0$) & $0.22 \pm 0.02$ & -- & $-0.01 \pm 0.05$ & $0.01 \pm 0.10$ & no punishment \\
\quad No memory & $0.95 \pm 0.01$ & -- & $0.00 \pm 0.00$ & $0.48 \pm 0.03$ & no punishment \\
\quad Uninformative memory & $0.053 \pm 0.003$ & -- & $0.00 \pm 0.00$ & $0.07 \pm 0.04$ & no punishment \\
\bottomrule
\end{tabular}}
\end{table}

The baseline reproduces \citet{calvano2020}: $\Delta=\BertBase$, and after a
deviation the rival cuts its price ($\BertBaseRival$) and the deviator loses
($\BertBaseGain$). The test sees punishment when it is there. The placebos
then behave as they should, and they also show why the price level is not
enough. Myopic learners, who by Proposition~\ref{prop:myopic} would play Nash
if their estimates were unbiased, settle at $\Delta=\BertMyopic$ and do not
punish; the departure from Nash is their learning bias. Memoryless learners
reach $\Delta=\BertNone$, \emph{higher} than the baseline, with no rival
reaction at all and a profitable deviation ($\BertNoneGain$): a price level
indistinguishable from collusion with nothing behind it, the ``spontaneous
coupling'' of \citet{banchio2023}. This differs from the modest gains that
\citet{calvano2020} report for memoryless algorithms, and the difference is
instructive. For $k=0$ they argue that there is no loss of generality in
setting $\delta=0$, and run memoryless learners with $\delta=0$, $\alpha=0.25$
and $\beta=10^{-4}$ (their online appendix, Section A4.1). With that
specification we reproduce their result ($\Delta=\BertNoneSpec$). Keeping
every parameter but setting $\delta=0.95$ gives $\Delta=\BertNoneSpecG$, again
with no punishment. Setting $\delta=0$ is innocuous for equilibrium, since a
memoryless strategy cannot condition on the past, but not for learning:
the discount factor changes the dynamics of the estimates, and with it the
outcome. An uninformative memory of the same size as
the real one collapses prices to $\Delta=\BertRandom$, so in this game what
memory contains matters, as punishment requires. With noisy profits the
baseline still gives $\Delta=\BertNoise$ and still punishes (myopic:
$\BertNoiseMyopic$). Counterfactual updates lower prices to
$\Delta=\BertCf$, as \citet{asker2022} find for synchronous learning, but what
remains is still enforced: the rival cuts its price after a deviation
($\BertCfRival$) and deviating loses ($\BertCfGain$). The learning rule
changes the level, not the nature, of collusion in this game.

The diagnostics therefore discriminate. They find punishment in the game
where the literature finds it, and in the same game they flag a
supracompetitive outcome without punishment when memory is removed. What
they find among traders is not a property of the tests.

\subsection{A second market: dealers under adverse selection}
\label{sec:dealers}

Q-learning market makers also sustain supracompetitive quotes
\citep{colliard2026}. We apply the same diagnostics to a stylised dealer
market. Two dealers post ask quotes on a 15-point grid spanning the
competitive and monopoly asks plus 10\%; each period one buyer arrives. With
probability $\pi=0.3$ the buyer is informed and buys only when the asset is
worth $v=+\sigma$ ($v=\pm\sigma$, $\sigma=0.5$); otherwise it is a liquidity
buyer with valuation $w\sim U[0,1]$ who buys if $w$ exceeds the best ask. The
dealer with the lower ask serves the order (ties split) and earns ask minus
$v$. The competitive ask is the zero-profit \citet{glosten1985} quote
($a^N=\QuoteAN$) and the monopoly ask maximises expected profit
($a^M=\QuoteAM$). Feedback is noisy, as the buyer's type and the asset value are
random, and the profit every ask would have earned against the same buyer is
known, so counterfactual updates apply exactly. Dealers remember both last
quotes, as in \citet{calvano2020}, and learn with the same parameters.

The lower panel of Table~\ref{tab:validation} reports the results. Dealers
quote well above the competitive ask: the best ask closes $\QuoteBase$ of the
gap to the monopoly ask and dealers earn $\QuoteBaseProfit$ of the monopoly
profit gap. Nothing in this outcome is enforced. After a forced undercut the
rival does not lower its quote ($\QuoteBaseRival$), and undercutting pays
($+\QuoteBaseGain$). Every placebo reproduces the markup: myopic dealers
quote even higher ($\QuoteMyopic$), as do dealers without memory
($\QuoteNone$) or with an uninformative one ($\QuoteRandom$). Counterfactual
updates, which remove the selection effect, bring the best ask most of the
way to the competitive level ($\QuoteCf$; without memory $\QuoteCfNone$).
This is the pattern of a learning bias, and it agrees with the reading of
\citet{colliard2026}, who trace dealers' markups to limited experimentation
and noisy feedback rather than to punishment. The protocol separates the two
markets cleanly: the same tests that find punishment in the pricing duopoly
find none among dealers, although both reach similar price levels.
